% !TeX program = pdflatex
\documentclass[11pt,a4paper]{amsart}

\newcommand{\notelanguage}{finnish}
% Shared preamble for course notes and exercise sets.

\usepackage[T1]{fontenc}
\usepackage{lmodern}
\usepackage[english]{babel}

\usepackage[a4paper,margin=30mm]{geometry}
\usepackage{microtype}
\usepackage{graphicx}
\usepackage{enumitem}

\usepackage{amsmath}
\usepackage{amssymb}
\usepackage{amsthm}
\usepackage{mathtools}
\usepackage{aliascnt}

\usepackage[hidelinks,pdfusetitle]{hyperref}

\numberwithin{equation}{section}
\setlist{itemsep=0.25em,topsep=0.5em}

\theoremstyle{plain}
\newtheorem{theorem}{Theorem}[section]
\newaliascnt{lemma}{theorem}
\newtheorem{lemma}[lemma]{Lemma}
\aliascntresetthe{lemma}
\newaliascnt{proposition}{theorem}
\newtheorem{proposition}[proposition]{Proposition}
\aliascntresetthe{proposition}
\newaliascnt{corollary}{theorem}
\newtheorem{corollary}[corollary]{Corollary}
\aliascntresetthe{corollary}

\theoremstyle{definition}
\newaliascnt{definition}{theorem}
\newtheorem{definition}[definition]{Definition}
\aliascntresetthe{definition}
\newaliascnt{example}{theorem}
\newtheorem{example}[example]{Example}
\aliascntresetthe{example}
\newaliascnt{problem}{theorem}
\newtheorem{problem}[problem]{Problem}
\aliascntresetthe{problem}

\theoremstyle{remark}
\newaliascnt{remark}{theorem}
\newtheorem{remark}[remark]{Remark}
\aliascntresetthe{remark}

\usepackage[nameinlink,noabbrev]{cleveref}

\crefname{theorem}{theorem}{theorems}
\crefname{lemma}{lemma}{lemmas}
\crefname{proposition}{proposition}{propositions}
\crefname{corollary}{corollary}{corollaries}
\crefname{definition}{definition}{definitions}
\crefname{example}{example}{examples}
\crefname{problem}{problem}{problems}
\crefname{remark}{remark}{remarks}

\DeclareMathOperator{\im}{im}
\DeclareMathOperator{\rank}{rank}
\DeclarePairedDelimiter{\abs}{\lvert}{\rvert}
\DeclarePairedDelimiter{\norm}{\lVert}{\rVert}
\DeclarePairedDelimiter{\set}{\lbrace}{\rbrace}

\usepackage{tikz}

\title{Analyysi I\\Harjoituksen 4 ratkaisut}
\author{}
\date{}

\begin{document}

\exerciseheading{Analyysi I}{Harjoituksen 4 ratkaisut}

\begin{exercise}
  Olkoon \(x \in \RR\). Määritellään jono \((x_n)\) asettamalla, että
  \[
    x_n = \left(1 + \frac{x}{n}\right)^n.
  \]
  Osoita, että jono \((x_n)\) on kasvava kun \(n > \abs{x}\).
\end{exercise}

\begin{solution}
  Olkoon \(n \in \NN\) ja \(n > \abs{x}\). Tällöin \(n + x > 0\), joten
  \(x_n > 0\). Riittää siis osoittaa, että \(\frac{x_{n+1}}{x_n} \geq 1\).
  Laskemalla saadaan
  \[
    \begin{aligned}
      \frac{x_{n+1}}{x_n} &= \frac{\left(1 + \frac{x}{n + 1}\right)^{n + 1}}
      {\left(1 + \frac{x}{n}\right)^n} \\
      &= \left(\frac{1 + \frac{x}{n + 1}}{1 + \frac{x}{n}}\right)^{n
      + 1}\left(1 + \frac{x}{n}\right) \\
      &= \left(\frac{n^2 + nx + n}{n^2 + nx + n + x}\right)^{n +
      1}\left(\frac{n + x}{n}\right) \\
      &= \left(1 - \frac{x}{n^2 + nx + n + x}\right)^{n +
      1}\left(\frac{n + x}{n}\right).
    \end{aligned}
  \]
  Asetetaan \(t = -\frac{x}{(n + 1)(n + x)}\). Koska \(n > \abs{x}\),
  saadaan
  \[
    1 + t = \frac{n(n + 1 + x)}{(n + 1)(n + x)} > 0,
  \]
  joten \(t > -1\) ja Bernoullin epäyhtälöä voidaan soveltaa.
  Lisäksi \(\frac{n + x}{n} > 0\), joten tällä luvulla kertominen säilyttää
  epäyhtälön suunnan. Näin saadaan arvio
  \[
    \begin{aligned}
      \frac{x_{n+1}}{x_n} &=
      \left(1 - \frac{x}{n^2 + nx + n + x}\right)^{n +
      1}\left(\frac{n + x}{n}\right) \\
      &\geq \left(1 - \frac{(n + 1)x}{n^2 + nx + n +
      x}\right)\left(\frac{n + x}{n}\right) \\
      &= \left(\frac{n^2 + n}{n^2 + nx + n + x}\right)\left(\frac{n +
      x}{n}\right) \\
      &= 1.
    \end{aligned}
  \]
  Olemme osoittaneet, että \(x_{n + 1} \geq x_{n}\) kaikilla \(n \in \NN\),
  joille \(n > \abs{x}\), eli lukujono \((x_n)\) on kasvava tällä
  indeksialueella.
\end{solution}

\begin{exercise}
  \leavevmode\par
\begin{enumerate}[label=\alph*), nosep]
  \item Osoita, että lukujonon \(\left(\left(1 - \frac{1}{n}\right)^n\right)\)
    raja-arvo on \(\frac{1}{e}\).
  \item Muodosta edellisten kohtien perusteella yläraja Neperin luvulle.
  \item Valitse \(n = 50, 500, 5000\), ja arvioi Neperin lukua ala- ja
    ylärajojen avulla.
\end{enumerate}
\end{exercise}

\begin{solution}
\leavevmode\par
\begin{enumerate}[label=\alph*), nosep]
\item Kun \(n \geq 2\), niin
  \[
    1 - \frac{1}{n} = \frac{n - 1}{n} = \left(1 + \frac{1}{n - 1}\right)^{-1}.
  \]
  Näin ollen
  \[
    \left(1 - \frac{1}{n}\right)^n = \frac{1}{
      \left(1 + \frac{1}{n - 1}\right)^{n - 1}\left(1 + \frac{1}{n - 1}\right)
    }.
  \]
  Koska
  \[
    \left(1 + \frac{1}{n}\right)^n \to e, \quad \text{kun} \quad n \to \infty,
  \]
  ja \(n - 1 \to \infty\), kun \(n \to \infty\), myös
  \[
    \left(1 + \frac{1}{n - 1}\right)^{n - 1} \to e, \quad \text{kun}
    \quad n \to \infty.
  \]
  Osoitetaan vielä, että \(1 + \frac{1}{n - 1} \to 1\), kun \(n \to \infty\).
  Olkoon \(\epsilon > 0\). Valitaan \(N \in \NN\) siten, että
  \(N > \max\{2, \frac{2}{\epsilon}\}\). Tällöin kaikilla \(n > N\) pätee
  \[
    \abs*{1 + \frac{1}{n - 1} - 1} = \frac{1}{n - 1} \leq \frac{2}{n}
    < \epsilon.
  \]
  Lisäksi vakiojono \(1\) suppenee kohti lukua \(1\).
  Koska \(e > 0\), nimittäjän raja-arvo \(e \cdot 1 = e\) on nollasta
  poikkeava. Tulon ja osamäärän raja-arvosääntöjen (lause 2.2.8) nojalla
  saadaan siis
  \[
    \left(1 - \frac{1}{n}\right)^n = \frac{1}{
      \left(1 + \frac{1}{n - 1}\right)^{n - 1}\left(1 + \frac{1}{n - 1}\right)
    } \to \frac{1}{e},
  \]
  kun \(n \to \infty\).
\item Tehtävän 1 mukaan lukujono
  \[
    \left(\left(1 - \frac{1}{n}\right)^n\right)
  \]
  on kasvava, kun \(n \geq 2\). Koska jonon termit ovat positiivisia,
  käänteislukujen muodostama jono
  \[
    \left(\left(1 - \frac{1}{n}\right)^{-n}\right)
  \]
on vähenevä. Kohdan a) mukaan
\[
  \lim_{n \to \infty} \left(1 - \frac{1}{n}\right)^n =
  \frac{1}{e},
\]
joten käänteisluvun raja-arvosäännön nojalla
\[
  \lim_{n \to \infty} \left(1 - \frac{1}{n}\right)^{-n} = e.
\]
Vähenevän ja suppenevan jonon jokainen termi on vähintään jonon
raja-arvon suuruinen. Näin saadaan yläraja
\[
  e \leq \left(1 - \frac{1}{n}\right)^{-n},
\]
kun \(n \geq 2\).
\item Tehtävän 1 ja tunnetun raja-arvon
\(\left(1 + \frac{1}{n}\right)^n \to e\) perusteella saadaan alaraja.
Yhdistämällä se kohdan b) ylärajaan saadaan
\[
\left(1 + \frac{1}{n}\right)^n \leq e
\leq \left(1 - \frac{1}{n}\right)^{-n}, \qquad n \geq 2.
\]
Lasketaan rajat arvoilla \(n = 50, 500, 5000\). Pyöristetään alarajat
alaspäin ja ylärajat ylöspäin kuuden desimaalin tarkkuuteen, jotta
epäyhtälöt säilyvät voimassa. Näin saadaan
\[
\begin{aligned}
  n = 50:   &\qquad 2{,}691588 \leq e \leq 2{,}745973, \\
  n = 500:  &\qquad 2{,}715568 \leq e \leq 2{,}721006, \\
  n = 5000: &\qquad 2{,}718010 \leq e \leq 2{,}718554.
\end{aligned}
\]
\end{enumerate}
\end{solution}

\begin{exercise}
Määritellään lukujono \((z_n)\) asettamalla
\(z_n = (-1)^n\left(2 + \frac{3}{n}\right)\). Osoita, että lukujono
hajaantuu. Sisältääkö lukujono suppenevan osajonon?
\end{exercise}

\begin{solution}
Olkoot \((x_n)\) ja \((y_n)\) lukujonon \((z_n)\) seuraavat osajonot:
\[
\begin{aligned}
x_n &= z_{2n} = 2 + \frac{3}{2n}, \\
y_n &= z_{2n + 1} = -2 - \frac{3}{2n + 1}.
\end{aligned}
\]
Olkoon \(\epsilon > 0\). Valitaan \(N \in \NN\) siten, että
\(N > \frac{3}{\epsilon}\). Tällöin kaikilla \(n > N\) pätee
\[
\abs{x_n - 2} = \frac{3}{2n} \leq \frac{3}{n} < \epsilon
\]
sekä
\[
\abs{y_n + 2} = \frac{3}{2n + 1} \leq \frac{3}{n} < \epsilon.
\]
Täten \(x_n \to 2\) ja \(y_n \to -2\), kun \(n \to \infty\).
Jos lukujono \((z_n)\) suppenisi, sen kaikki osajonot suppenisivat samaan
raja-arvoon. Koska \(2 \neq -2\), lukujono \((z_n)\) hajaantuu.
Se sisältää kuitenkin suppenevat osajonot \((x_n)\) ja \((y_n)\).
\end{solution}

\begin{exercise}
Olkoon \((z_k)\) vähenevä, alhaalta rajoitettu lukujono. Osoita, että
\((z_k)\) suppenee, ja että
\(\lim_{k \to \infty} z_k = \inf\{z_k : k \in \NN\}\).
\end{exercise}

\begin{solution}
Olkoon \((z_n)\) vähenevä ja alhaalta rajoitettu lukujono. Koska sen arvojen
joukko on epätyhjä ja alhaalta rajoitettu, on olemassa
\[
p = \inf{\set{z_n \colon n \in \NN}} \in \RR.
\]
Olkoon \(\epsilon > 0\). Lauseen 1.2.8 nojalla on olemassa \(N \in \NN\),
jolle pätee
\[
p \leq z_N < p + \epsilon.
\]
Koska \((z_n)\) on vähenevä,
\[
p \leq z_n \leq z_N < p + \epsilon,
\]
kun \(n > N\). Toisin sanoen
\[
\abs{z_n - p} = z_n - p < \epsilon,
\]
kun \(n > N\). Täten \(\lim_{n \to \infty} z_n = p\).
\end{solution}

\begin{exercise}
Tässä tehtävässä osoitetaan, että jokaisella lukujonolla on monotoninen
osajono. Täydennä seuraava todistus loppuun: Olkoon \((x_n)\) lukujono.
Sanomme, että \(x_n\) on lukujonon huippu, jos \(x_n \geq x_m\) kaikilla
\(m \geq n\). Jos huippuja on ääretön määrä \(x_{n_k}\), missä
\(n_k < n_{k+1}\) kun \(k \in \NN\), niin ne muodostavat laskevan
lukujonon. Jos lukujonon \((x_n)\) huippuja on äärellinen määrä, niin
viimeisen huipun jälkeen olevista alkioista löydämme kasvavan lukujonon.
\end{exercise}

\begin{solution}
Olkoon
\[
H = \set{m \in \NN \colon x_m \geq x_n \text{ kaikilla } n \geq m}.
\]
Joukko \(H\) koostuu siis jonon \((x_n)\) huippukohtien indekseistä.
Jaetaan todistus kahteen tapaukseen.
\begin{enumerate}
\item Oletetaan, että \(H\) on ääretön. Luetellaan sen alkiot kasvavassa
järjestyksessä \(m_1 < m_2 < m_3 < \cdots\) ja tarkastellaan osajonoa
\((x_{m_k})\). Koska \(m_k \in H\), huippukohdan määritelmän mukaan
\(x_j \leq x_{m_k}\) kaikilla \(j \geq m_k\). Erityisesti
\(m_{k+1} > m_k\), joten
\[
x_{m_{k+1}} \leq x_{m_k} \qquad \text{kaikilla } k \in \NN.
\]
Siis \((x_{m_k})\) on lukujonon \((x_n)\) vähenevä osajono.
\item Oletetaan, että \(H\) on äärellinen. Jos \(H = \varnothing\),
asetetaan \(k_1 = 1\). Muussa tapauksessa asetetaan \(k_1 = \max H + 1\).
Tällöin mikään indeksi \(j \geq k_1\) ei kuulu joukkoon \(H\).

Muodostetaan osajono rekursiivisesti. Kun \(k_n \geq k_1\) on valittu,
pätee \(k_n \notin H\). Huippukohdan määritelmän mukaan on siis olemassa
indeksi \(k_{n+1} > k_n\), jolle
\[
x_{k_{n+1}} > x_{k_n}.
\]
Koska myös \(k_{n+1} \geq k_1\), konstruktiota voidaan jatkaa kaikilla
\(n \in \NN\). Näin saadaan lukujonon \((x_n)\) aidosti kasvava
osajono \((x_{k_n})\).
\end{enumerate}
Molemmissa tapauksissa lukujonolla \((x_n)\) on siis monotoninen osajono.
\end{solution}

\begin{exercise}
Olkoon \((x_n)\) lukujono, joka toteuttaa ehdon
\(\abs{x_n - x_{n+1}} < 5 \cdot 2^{-n}\) kaikilla \(n \in \NN\).
Osoita Cauchyn yleisen suppenemisehdon avulla, että lukujono \((x_n)\)
suppenee.

Tehtävässä saa käyttää tietoa
\[
1 + x + \ldots + x^n = \frac{1 - x^{n+1}}{1 - x},
\]
kun \(n \in \NN\) ja \(x \in \RR\), \(x \neq 1\)
(kurssikirjan Lemma 1.4.10).
\end{exercise}

\begin{solution}
Olkoot \(n,k \in \NN\), \(n,k \geq 1\). Teleskooppisumman,
kolmioepäyhtälön ja tehtävän oletuksen avulla saadaan
\[
\begin{aligned}
\abs{x_n - x_{n + k}} &= \abs*{\sum_{j = 0}^{k - 1}
\left(x_{n + j} - x_{n + j + 1}\right)} \\
&\leq \sum_{j = 0}^{k - 1} \abs{x_{n + j} - x_{n + j + 1}} \\
&< 5\sum_{j = 0}^{k - 1} 2^{-(n + j)}
= 5\cdot2^{-n}\sum_{j = 0}^{k - 1} \left(\frac12\right)^j.
\end{aligned}
\]
Tehtävässä annetun geometrisen summan kaavan (lemma 1.4.10) nojalla
\[
\sum_{j = 0}^{k - 1} \left(\frac12\right)^j
= \frac{1-2^{-k}}{1-\frac12}
= 2(1-2^{-k}) < 2.
\]
Saadaan siis indeksistä \(k\) riippumaton arvio
\[
\abs{x_n-x_{n+k}} < 5\cdot2^{-n}\cdot2 = \frac{10}{2^n}.
\]
Olkoon \(\epsilon > 0\). Valitaan \(N \in \NN\) siten, että
\(N > \frac{10}{\epsilon}\). Koska \(2^n > n\) kaikilla \(n \geq 1\),
kaikilla \(n > N\) ja \(k \geq 1\) pätee
\[
\abs{x_n - x_{n + k}} < \frac{10}{2^n} < \frac{10}{n} < \epsilon.
\]
Kun \(k=0\), pätee \(\abs{x_n-x_n}=0<\epsilon\). Näin ollen
\(\abs{x_n-x_{n+k}}<\epsilon\) kaikilla \(n>N\) ja \(k\in\NN_0\).
Cauchyn yleisen suppenemisehdon (lause 2.4.4) nojalla lukujono
\((x_n)\) suppenee.
\end{solution}

\end{document}
